% Propietario conservado: /Users/ruben/Documents/New project/output/EXTREMA_MEDIA_RAZON_RECIPROCIDAD_ES_EN_20260918/ARTICULO/ES/source/libro/06_sintesis.tex
\chapter{Unité dimensionnelle, symétrie relative et récupération de la structure}
\label{x_reciprocidad:lib06:sintesis}

L’extrême et moyenne raison holographique articule la conservation évolutive
de l’unité avec ses réalisations arithmétiques et dimensionnelles. La
composition développée permet de préciser ce principe au moyen de
quantités conservées, de domaines et d’inverses. L’état enrichi
engendré par APP--TRIT--TPK précède les lectures numériques et
géométriques ; ses opérations conservent l’orientation, la retenue,
l’incidence et la mémoire que chaque réalisation utilise.

\section{Conservation arithmétique, unitarité et décomposition isométrique}
\label{x_reciprocidad:lib06:tres-conservaciones}

Le raffinement arithmétique conserve l'unité normalisée :
\begin{equation}
 (x,y,c)\longmapsto(Bx-c,By+c),\qquad
 \frac{Bx-c}{B(x+y)}+\frac{By+c}{B(x+y)}=1.
 \label{x_reciprocidad:lib06:unidad}
\end{equation}
L'inverse par résidus récupère la retenue canonique et la paire précédente. La composition des blocs conserve les contributions des événements avec leur échelle, et les divisions successives récupèrent toute histoire finie. La lecture limite est contrôlée au moyen de cylindres décroissants ; les marques de frontière maintiennent la distinction entre les histoires qui partagent une évaluation scalaire.

Le relèvement sur les étiquettes transforme cette bijection en un
unitaire \(J\). Il conserve la norme des amplitudes et transporte
les lecteurs de fraction et de retenue au moyen de
\(J^*M_hJ=M_{h\circ\ell}\). La quantité \(x^2+y^2\) et la norme
de \(|x,y,c\rangle\) appartiennent à des réalisations distinctes. La
preuve d’unitarité établit le lien précis entre la
réversibilité arithmétique et la conservation des amplitudes.

La troisième conservation correspond à deux réalisations isométriques :
\begin{equation}
 (a+1)Q_-^*Q_-+aQ_+^*Q_+
 =(2a+1)(a^2+a+1)I.
 \label{x_reciprocidad:lib06:balance}
\end{equation}
L’annulation des termes mixtes permet la distribution
bilatérale et sa reconstruction de norme un. La loi vaut pour les
transports du domaine spécifié dans la
\cref{x_reciprocidad:lib04:teo-balance}. Sa spécialisation au cycle ternaire
apporte des propriétés supplémentaires de norme et d’indice ; la loi générale
conserve sa validité pour des isométries dépourvues de cet ordre fini.

La chaîne formée par les trois applications est effective : la bijection produit \(J\) ; une permutation préalablement déclarée produit le second transport \(JR\) ; les deux transports satisfont le bilan, et leurs canaux récupèrent l'entrée par \eqref{x_reciprocidad:lib05:inversa-bilateral}. L'exemple de la \cref{x_reciprocidad:lib05:ejemplo} montre la procédure complète avec \((73,27,1)\), \((729,271)\) et \(1001\) étiquettes admissibles. Les nombres de cette réalisation ne sont pas introduits comme substituts à la structure qui engendre le registre \(K\).

\section{Mémoire complète et contribution logique à la limite de Landauer}
\label{x_reciprocidad:landauer:puente-archivo}

La contribution logique nulle de Landauer, établie dans
\eqref{x_reciprocidad:v:ant:eq:ii-kb-aut-reversible}, s’applique aux transports
récupérables construits dans cet article. La bijection de
\cref{x_reciprocidad:lib05:teo-biyeccion} conserve la retenue par
\eqref{x_reciprocidad:lib05:inversa-etiquetas} ; son prolongement conserve la route
au moyen de \cref{x_reciprocidad:lib05:teo-rutas}. Les archives de
\cref{x_reciprocidad:lib04:teo-archivo} conservent les amplitudes et les corrélations
nécessaires pour inverser le transport.

\begin{corollary}[Récupération et conservation spectrale de l’archive]
\label{x_reciprocidad:landauer:cor-archivo}
Soit \(Z:H\to\mathcal K\) l’une des archives complètes
\(Z_N,Z_\infty\) de \cref{x_reciprocidad:lib04:teo-archivo}, ou le relèvement
unitaire \(J_n\) de \cref{x_reciprocidad:lib05:teo-rutas}, dans son domaine respectif.
Pour un opérateur de densité de rang fini \(\rho\) sur \(H\),
la sortie \(\rho_Z=Z\rho Z^*\) satisfait
\begin{align}
 \operatorname{Tr}\rho_Z&=1,& Z^*\rho_ZZ&=\rho,\\
 s(\rho_Z)&=s(\rho),&
 s(\rho)&=-\operatorname{Tr}(\rho\log\rho),
 \label{x_reciprocidad:landauer:espectro}
\end{align}
avec \(0\log0=0\). Si \(X\) est une variable aléatoire sur les routes
finies de \cref{x_reciprocidad:lib05:teo-rutas}, alors
\begin{equation}
 \mathsf H(X\mid\Phi_n(X))=0.
 \label{x_reciprocidad:landauer:recuperacion-ruta}
\end{equation}
\end{corollary}
\begin{proof}
Écrivons la décomposition spectrale finie \(\rho=\sum_jp_j|v_j\rangle\langle v_j|\), avec les \(v_j\) orthonormaux, \(p_j>0\) et \(\sum_jp_j=1\). L'égalité \(Z^*Z=I\) implique que les \(Zv_j\) sont orthonormaux. Par conséquent, \(\rho_Z=\sum_jp_j|Zv_j\rangle\langle Zv_j|\) possède les mêmes valeurs propres non nulles, ce qui prouve la trace et l'entropie. Multiplier par \(Z^*\) et \(Z\) récupère \(\rho\). Enfin, \(\Phi_n^{-1}\) récupère chaque chemin à partir de sa paire finale ; la distribution conditionnelle est concentrée sur un seul chemin et son entropie est nulle.
\end{proof}

La conservation correspond à l’archive complète. Une lecture visible
\(q\), accompagnée de mémoire \(M\), retrouve l’état sous
l’hypothèse de \eqref{x_reciprocidad:v:ant:eq:ii-kb-aut-fibra} ; l’information non
publiée par \(q\) reste dans
\(\mathsf H(M(X)\mid q(X))\). Écarter un complément ou un
terminal nécessaire définit une autre opération, dont la récupérabilité est
déterminée sur l’information effectivement retenue. L’égalité
spectrale précédente conserve la distinction entre l’entropie de l’état
complet, l’entropie d’une réduction partielle et la chaleur transférée.

Dans le modèle thermique de la section
\ref{x_reciprocidad:v:ant:sec:ii-kb-aut-desarrollo}, avec états logiques dégénérés,
environnement isotherme et comptabilité des registres auxiliaires, le terme
associé à l’effacement est
\(k_B\Theta\,\mathsf H(X\mid F(X))\), pour une application logique
déterministe \(F\). Sa contribution est nulle pour les transports
récupérables précédents ; la réinitialisation d’un TRIT uniforme satisfait
\(Q_{\rm reset}\geq k_B\Theta\log3\). La préparation, le contrôle,
la vitesse finie et le couplage à l’environnement conservent leurs
bilans physiques propres. La chaleur totale d’une implémentation
requiert de composer ces contributions avec l’opération logique.

\section{Conservation des observables et symétrie relative}
\label{x_reciprocidad:lib06:simetria}

Le bilan de la norme est le cas de l’identité d’une loi opératoire
plus informative. Lorsque les transports sont unitaires, de relatif
\(R\), le lecteur effectif est
\begin{equation}
 \mathcal T_a(G)=\frac{a(a+1)G+R^*GR}{a(a+1)+1}.
 \label{x_reciprocidad:lib06:lector}
\end{equation}
L'observable reste fixe exactement lorsque \([G,R]=0\). Dans le cas canonique,
\begin{equation}
 \mathcal T_8(G)-G=\frac{R^*GR-G}{73}.
 \label{x_reciprocidad:lib06:defecto-lector}
\end{equation}
Cette identité détermine aussi bien la conservation que la variation : le défaut mesure la différence entre deux réalisations du même lecteur, pondérée par la répartition \(72+1\).

La mémoire complète transporte en outre l’observable au moyen de
\(A\mapsto ZAZ^*\) dans \(\operatorname{im}Z\). Les produits,
adjoints et appariements sont conservés dans cette image. La lecture
diagonale des canaux et l’observable transporté sont des opérations
différentes ; l’exemple de la retenue montre la différence entre
\(74/73\) pour la première et la valeur propre un pour le second.
L’information demeure récupérable même si une lecture diagonale
particulière change de valeur.

Les charges variationnelles de la \cref{x_reciprocidad:nuc11:noether} conservent leur
fondement dans une action et un entrelacement explicites. Pour
\(A_{n+1}U_n=U_nA_n\), l’action de transport donne la charge
\(p_n^*A_nq_n\). Sa conservation et l’identité de norme sont des
résultats compatibles, sous des hypothèses différentes. De même,
la conservation d’une forme alternée ou lorentzienne dans
\eqref{x_reciprocidad:lib04:balance-formas} maintient son sens géométrique ;
les estimations de bruit et de contraction utilisent la positivité
hilbertienne déclarée dans leurs preuves.

\section{Fonction du registre dans la transduction}
\label{x_reciprocidad:lib06:registro}

La constante holographique de couplage arithmético-géométrique possède une fonction définie par ses compositions. Le registre \(K\) est obtenu préalablement à partir de l'extraction orientée et de Hadamard. Sa lecture \(\kappa\) est réversible avec une base, une longueur et une origine fixées. Son repère \(O_K\) identifie des réponses cycliques car sa matrice de Gram possède une borne inférieure strictement positive. L'incidence et la direction géométrique utilisent leurs lecteurs déclarés sur le même registre.

L’isométrie d’archive conserve le gramien de ce repère, de sorte qu’elle
maintient les bornes de récupération à toute profondeur. Le
registre complet admet une inversion de norme un ; les deux mémoires
partielles conservent leur facteur \(9/4\) et leur noyau uniforme. La
charge supplémentaire distingue la restitution du registre absolu de
celle de sa composante centrée. Le décodage entier restitue
ensuite la sélection incidentielle même à un seuil situé
exactement sur une coordonnée.

Cette capacité est plus spécifique que la conservation d’une grandeur totale :
elle permet de retrouver des opérations et des relations qui agissent sur un
objet engendré. Ses garanties reposent sur le Gram, sur les
distances de l’image entière et sur les inverses démontrés, et non
sur la conservation par définition d’une autre copie de l’entrée. En même temps,
le nombre fini de fenêtres de \(K\) reste distinct de
l’histoire enrichie complète ; les autres registres et marques
qui l’accompagnent maintiennent leurs transports propres.

L'orbite et ses transports isométriques conservent la norme \(\nu=\|K\|\). Normaliser ses colonnes à une norme égale à un exige de conserver \(\nu\) pour inverser cette opération. La normalisation polaire utilise la paire \((F_K,G_K)\). La formule \(O_K=F_KG_K^{1/2}\) maintient l'orientation et les amplitudes. Le lecteur positionnel doit accompagner le changement de carte, car la norme orbitale commune ne rend pas \(\kappa\) invariant sous une renumérotation isolée de ses composantes.

\section{Profondeur, dimensions et réalisations physiques}
\label{x_reciprocidad:lib06:profundidad}

La politique bilatérale qui continue un canal et archive l’autre
a un épuisement uniforme du terminal. Sa limite est une
isométrie vers la somme de toutes les mémoires, et son inverse est
la série adjointe convergente de
\eqref{x_reciprocidad:lib04:recuperacion-infinita}. Un nombre fini de
compléments possède deux reconstructions différenciées :
l’adjoint tronqué conserve un biais explicite et l’inverse des
moindres carrés l’élimine avec l’amplification du bruit
calculée. Les évolutions précédentes à terminal positif
conservent cette composante dans l’archive complète.

La dimension de l’espace des étiquettes, le rang d’un projecteur et
la dimension d’une réalisation physique sont reliés par
leurs applications concrètes. L’unitaire à \(1001\) étiquettes et le
registre de douze fenêtres sont deux réalisations explicites de
domaines distincts. La conservation extérieure du développement
précédent prolonge les produits de Gram à chaque degré fini,
en incluant les contributions mixtes entre terminal et mémoires.
Les réalisations exceptionnelles et physiques restent liées
à leurs représentations et opérateurs déjà construits ; les rangs
numériques ne remplacent pas à eux seuls ce lien.

\Needspace{18\baselineskip}
\section{Conclusions sur la conservation évolutive de l'unité}
Les actions physiques et le protocole gaussien du manuscrit conservent les couplages, les formes, les unités et les préparations qui interviennent dans leurs preuves. Le bilan bilatéral ajouté ici élargit la famille des transports conservatifs et de leurs lectures ; son application à une réalisation physique s'établit en composant les opérateurs de cette réalisation, sans les remplacer par une coïncidence de normalisation.

L'unité du développement réside donc dans une séquence démontrée : génération et complétion, raffinement avec mémoire, réalisation unitaire, symétrie relative du lecteur et récupération de la structure. L'extrême et moyenne raison holographique reçoit ainsi un exposé dans lequel changement d'échelle, incidence, conservation et transduction possèdent des opérations explicites et maintiennent conjointement l'information nécessaire pour inverser leurs lectures.
